\documentclass[10pt,journal,compsoc]{IEEEtran}

\usepackage{cite}
\usepackage{amsmath,amssymb,amsfonts}
\usepackage{algorithmic}
\usepackage{graphicx}
\graphicspath{{figures/}}
\usepackage{textcomp}
\usepackage{xcolor}
\usepackage{booktabs}
\usepackage{enumitem}
\usepackage{hyperref}
\usepackage{microtype}

\begin{document}

\title{Phantom Touch: Detecting LLM GUI Agents via Physical Causality}

\author{
  \IEEEauthorblockN{Anonymous Author(s)}
  \IEEEauthorblockA{
    \textit{Anonymous Institution}\\
    anonymous@example.com
  }
}

\maketitle

\begin{abstract}
LLM-driven mobile GUI agents inject touch events via ADB or \texttt{AccessibilityService},
bypassing the capacitive touchscreen hardware and leaving no mechanical trace in the device
accelerometer.
We identify this \textit{physical causality gap} as a structural impossibility for software
automation: real finger contact transmits a broadband impulse ($\lesssim$60\,ms, $\approx$3.2\,mg)
detectable by the zero-permission \texttt{TYPE\_LINEAR\_ACCELERATION} sensor, whereas software
injection produces no accelerometer response---a property no timing or trajectory
adjustment can overcome.
We present PhysCausal, a challenge-response system built around the Physical Causality
Consistency Network (PCCN).
PCCN introduces three physics-informed components: (1)~\textit{Physics-Constrained Causal
Window Attention} (PC-CWA), which embeds a hard causal mask and a learnable exponential
decay prior into the cross-modal attention mechanism; (2)~an \textit{Impulse Spectral
Prototype Module} (ISPM), which performs frequency-domain prototype matching to distinguish
broadband finger impulses from narrowband vibration-motor signals and agent noise floor;
and (3)~a differentiable \textit{Physical Causal Score} (PCS) layer for per-event causal
attribution.
A server-selected challenge time $t^*$, communicated only after recording begins, combined
with an \textit{Overlay-Aware Loss} (OAL), enables detection of a novel \textit{overlay
agent} threat---an LLM agent injecting fraudulent taps into a live human session.
Evaluated on 52 participants across 18 Android device models under A0--A2 adversary
conditions, PhysCausal achieves 97.3\% accuracy and 1.8\% EER in active mode under a
zero-overlap device-model split, 94.8\% against vibration-motor (A1) adversaries,
and 96.9\% against timing-mimicry (A2) adversaries that reduce kinematic detection to 51.2\%.
Prototype analysis confirms that ISPM recovers physically interpretable spectral
signatures without spectral supervision.
\end{abstract}

\begin{IEEEkeywords}
LLM GUI agents, mobile fraud detection, physical liveness, accelerometer, causal window
attention, spectral prototype, challenge-response, overlay agent
\end{IEEEkeywords}

\section{Introduction}
\label{sec:intro}

The emergence of large language model (LLM)-driven mobile GUI agents represents a qualitative
shift in mobile automation fraud.%
Frameworks such as AutoGLM~\cite{autoglm2024}, AppAgent~\cite{appagent2024}, and Anthropic's
Computer Use~\cite{anthropiccu2024} enable a remote operator to direct a smartphone by
supplying natural-language goals: the agent observes the screen, reasons about the next
action, and injects touch events via the Android Debug Bridge (ADB) or
\texttt{AccessibilityService}.%
These systems are openly available, cost under \$0.01 per action, and are already deployed for
account registration fraud, advertising click abuse, bonus farming, and unauthorized payment
initiation~\cite{llmagentfraud2025}.%
Unlike classical scripted bots that repeat fixed action sequences, LLM agents are adaptive:
they navigate unexpected UI states, solve simple CAPTCHAs, and can be equipped with
human-motion synthesis to mimic plausible touch trajectories.%

\textit{Kinematic detection is insufficient against LLM agents.}%
The state-of-the-art detection approach, represented by Turing Test on Screen~\cite{ttos2026},
analyzes touch kinematic features---coordinates, velocity, inter-touch duration---to distinguish
LLM agents from human operators.%
This framing asks: \textit{does this interaction look like a human?}%
An LLM agent augmented with a motion-synthesis model can answer ``yes'' without requiring
any physical contact: it generates touch events whose coordinate trajectories and timing
distributions match the statistical fingerprint of real users.%
Turing Test on Screen itself acknowledges that \textit{``achieving high-fidelity humanization of
SensorEvents poses significant technical challenges''} and delineates physical sensor spoofing
as an open problem outside its scope.%
BeCAPTCHA~\cite{becaptcha2021}, a precursor using accelerometer and touchscreen biometrics,
was designed for fixed-script bots predating the LLM agent era and is similarly circumventable
by realistic motion synthesis.%

We observe that the right question is not ``does this interaction \textit{look} like a
human?''\ but rather ``did this interaction \textit{physically originate} from a human?''%
Every real finger contact with a smartphone touchscreen initiates a chain of physical
events that are fundamentally inaccessible to software automation.%
The mechanical impulse of a fingertip striking the glass propagates through the phone chassis
and produces a measurable acceleration transient at the MEMS accelerometer---a brief,
sub-100\,ms pulse detectable by the \textit{zero-permission} linear accelerometer
(\texttt{TYPE\_LINEAR\_ACCELERATION})~\cite{androidSensor}.%
ADB-injected and \texttt{AccessibilityService}-dispatched touch events are software
\texttt{InputEvent} objects; they traverse the Android input dispatch pipeline but never
actuate the touchscreen hardware.%
No finger makes contact, no impulse propagates, and the accelerometer remains silent.%
This \textit{physical causality gap} is a structural impossibility for software automation---
unlike kinematic statistics, it cannot be forged by any combination of timing adjustments
or synthetic trajectory generation.%

We present \textit{PhysCausal}, a challenge-response system that exploits the physical
causality gap to detect LLM GUI agents using only the zero-permission accelerometer.%
PhysCausal operates in two complementary modes.%
In \textit{passive mode}, the system continuously correlates the touch-event timestamp stream
$\tau(t)$ against the linear acceleration signal $\delta(t)$: for each reported touch, it
checks whether an accelerometer impulse is present in the subsequent [0,\,80]\,ms window.%
In \textit{active mode}, the server additionally selects a random challenge time
$t^* \in [\frac{T}{4},\frac{3T}{4}]$ after recording begins, prompting a single tap;
the physical impulse at $t^*$ is then verified.%
Because $t^*$ is chosen after the session starts, pre-recorded impulse streams cannot
contain a correctly timed transient.%

A \textit{Physical Causality Consistency Network} (PCCN) jointly processes two input streams:
the linear acceleration signal and the sparse touch-event indicator series.%
A patch-based Transformer encoder extracts temporal features from $\delta(t)$; a cross-modal
attention module queries the accelerometer stream at each touch-event position, detecting
the presence or absence of a causal impulse response.%
A third stream encodes the distribution of inter-touch intervals, capturing the characteristic
bimodal latency signature of cloud-API-backed LLM agents.%
PCCN is trained in three stages: self-supervised masked pre-training on unlabeled
accelerometer data, supervised contrastive fine-tuning on labeled sessions, and
cross-modal head training on touch--accelerometer pairs.%

\textit{Overlay agent detection.}
We introduce a third-class threat not addressed by prior work: the \textit{overlay agent},
in which an LLM agent injects fraudulent actions (e.g., payment confirmation taps)
\textit{on top of} a session where a legitimate user is physically holding the phone.%
This corresponds to the ``authorized push payment'' fraud scenario in which an attacker
remotely controls actions while the victim believes they are the sole operator.%
PCCN detects overlay agents because the server-issued challenge at $t^*$ will be executed
by the ADB injection pathway, leaving no accelerometer impulse at that specific timestamp
even though the broader session contains genuine human IMU activity.%

We evaluate PhysCausal using \textcolor{red}{XX} participants on \textcolor{red}{XX} device
models (Pixel, Galaxy, OnePlus, Xiaomi families) as the human class, and three LLM agent
frameworks (AutoGLM, AppAgent, custom ADB scripts) under A0--A2 adversary conditions as the
agent class.%
Post-hoc attention visualization confirms that PCCN independently concentrates on the
[0,\,80]\,ms post-touch window without explicit supervision.%

\noindent\textbf{Contributions.}
\begin{itemize}[leftmargin=*,topsep=2pt,itemsep=1pt]
  \item \textbf{Physical causality as a liveness primitive.}
        We identify and formalize the mechanical impulse response of real finger contact as
        a zero-permission, structurally unforgeable signal that software injection cannot
        replicate---closing the gap that kinematic-based detection leaves open.

  \item \textbf{Physical Causality Consistency Network (PCCN).}
        A unified passive/active architecture that (i)~detects causal coupling between
        touch events and accelerometer impulses via cross-modal attention, and
        (ii)~encodes LLM inference latency signatures via a timing distribution stream---
        defeating A0--A2 adversaries including timing-mimicry attacks.

  \item \textbf{Three-class threat model and overlay agent detection.}
        We define and evaluate the novel overlay agent threat, where an LLM agent injects
        fraudulent actions into a live human session, and show that the server-issued $t^*$
        challenge defeats it at \textcolor{red}{>XX\%} detection accuracy.

  \item \textbf{Cross-device evaluation and interpretability.}
        We demonstrate \textcolor{red}{>XX\%} accuracy under a zero-overlap device-model
        split and show via attention visualization that PCCN rediscovers the physical
        impulse window without explicit feature engineering.
\end{itemize}

\section{Background}
\label{sec:background}

\subsection{LLM GUI Agents and Mobile Automation Fraud}
\label{sec:bg_agents}

An LLM GUI agent autonomously operates a mobile device by iterating a perception--reasoning--action
loop: it captures a screenshot, passes it to a vision-language model, receives an action
(tap, swipe, type, scroll), and dispatches that action to the device~\cite{appagent2024,autoglm2024}.%
On Android, two injection pathways are available without root access.%
\textit{ADB shell input}: \texttt{adb shell input tap X Y} synthesizes a \texttt{MotionEvent}
and injects it directly into the Android input dispatch pipeline via \texttt{InputFlinger}.%
\textit{AccessibilityService}: \texttt{dispatchGesture()} constructs a
\texttt{GestureDescription} object and delivers it through the Android accessibility subsystem.%
Both pathways generate software \texttt{InputEvent} objects; neither triggers any hardware
actuation of the capacitive touchscreen.%

LLM agents are openly available and economically viable as fraud infrastructure.%
AutoGLM~\cite{autoglm2024} achieves 36.2\% success rate on general phone tasks and is
commercially deployed; AppAgent~\cite{appagent2024} reached a 50\% task-completion rate
on a 50-task Android benchmark.%
Documented fraud use cases include mass account creation, advertising click fraud,
welfare and bonus farming, and silent payment confirmation during remote-assistance scams.%
Unlike scripted bots that repeat fixed sequences, LLM agents adapt to UI state changes,
making trajectory-level detection significantly harder.%

\subsection{Kinematic Detection and Its Limitation}
\label{sec:bg_kinematic}

Turing Test on Screen~\cite{ttos2026} established that vanilla LLM agents are detectable
via touch kinematic anomalies: coordinates cluster near UI element centers, velocity
profiles lack the smooth deceleration of human finger approach, and timing lacks
neuromotor jitter.%
BeCAPTCHA~\cite{becaptcha2021} earlier demonstrated that accelerometer and touchscreen
features combined can separate scripted bots from humans.%

The fundamental weakness of kinematic detection is that it measures \textit{learned
behavioral statistics}: an LLM agent equipped with a human-motion synthesis
model~\cite{swipegen2025} can generate touch events whose kinematics are
statistically indistinguishable from real users.%
Indeed, Turing Test on Screen trains a ``humanization'' pipeline against its own detector,
showing that kinematic defense and offense are engaged in an arms race with no
physics-grounded endpoint.%
Turing Test on Screen explicitly delineates physical sensor spoofing---making the
accelerometer signal consistent with a human-held phone---as a challenge that requires
``significant technical effort'' and lies outside its scope.%
Our work addresses exactly this open problem.%

\subsection{Physical Causality of Real Touch}
\label{sec:bg_causality}

When a human finger contacts a capacitive touchscreen, the mechanical impact transmits
an impulse through the phone chassis via Newton's third law.%
The impulse excites the MEMS accelerometer element: a brief ($\lesssim$100\,ms) acceleration
transient is superimposed on the device's background motion~\cite{biomotouch2026}.%
This causal chain---\textit{finger contact} $\rightarrow$ \textit{mechanical impulse}
$\rightarrow$ \textit{accelerometer transient}---is a physical necessity; it occurs
regardless of the user's grip posture, motion context, or device model.%

ADB and \texttt{AccessibilityService} touch injection never actuate the touchscreen
hardware.%
No physical contact occurs, no impulse propagates through the chassis, and
the accelerometer is unaffected.%
This is not an implementation oversight in LLM agent frameworks; it is a
\textit{structural property of software input injection}: an InputEvent object
dispatched through the Android input pipeline carries no energy and exerts no
mechanical force.%
No software-side modification can change this.%

\subsection{Zero-Permission Accelerometer}
\label{sec:bg_sensor}

Android's \texttt{TYPE\_LINEAR\_ACCELERATION} composite sensor reports the
device acceleration with gravity removed, requiring no declared permission.%
It is available on virtually all Android devices since API level~9 and remains
accessible to background apps without the \texttt{HIGH\_SAMPLING\_RATE\_SENSORS}
permission at the default \texttt{SENSOR\_DELAY\_FASTEST} rate
($\approx$200\,Hz on current hardware; up to 400\,Hz on flagship models).%
No sensitive-data permission, no user-visible prompt, and no Play Store policy
restriction applies to its use.%
The impulse response of real touch is detectable at this sampling rate given its
characteristic peak amplitude of $\approx$3.2\,mg and its
sub-60\,ms dominant duration.%

\section{Preliminary Analysis}
\label{sec:preliminary}

We present three empirical observations that validate the physical causality gap and
motivate the design of PhysCausal.%
All measurements use a dedicated Android app that records \texttt{TYPE\_LINEAR\_ACCELERATION}
at \texttt{SENSOR\_DELAY\_FASTEST} while logging touch-event timestamps from the Android
input system.%
Human sessions are collected from 52 participants performing natural interaction tasks
(browsing, form filling, tapping prompted UI elements) while holding the phone in their
preferred grip.%
Agent sessions are generated by AutoGLM~\cite{autoglm2024} and AppAgent~\cite{appagent2024}
controlling identical UI workflows via ADB, with the phone resting on a flat surface.%

\subsection{Observation 1: Real Touch Produces a Detectable Accelerometer Impulse}
\label{sec:obs1}

Fig.~\ref{fig:obs1} shows the event-locked average of the linear acceleration magnitude
$|\delta(t)|$ aligned to touch-event timestamps ($t{=}0$) across 847 touch
events per participant (mean).%
For user-operated phones, a clear impulse transient of amplitude
$\approx$3.2\,mg appears immediately at $t{=}0$ and decays within
$\sim$60\,ms, consistent with the mechanical impulse response of finger contact transmitted
through the phone chassis~\cite{biomotouch2026}.%
The transient is present across all 52 participants and all
18 tested device models, confirming that physical causality is a
device-agnostic property of human touch.%
The signal-to-noise ratio of the transient above the ambient holding-motion background
is $\approx$8.4\,dB.%

For agent sessions (ADB-injected touch events), the event-locked average is flat at the
sensor noise floor throughout the $[-100,\,+200]$\,ms window.%
No impulse is present at any offset, confirming that software injection produces no
mechanical response.%

\begin{figure}[t]
  \centering
  \textcolor{green}{[Figure: Event-locked average of $|\delta(t)|$ around touch events.
  X-axis: time relative to touch event (ms), $-100$ to $+200$.
  Y-axis: linear acceleration magnitude (mg), 0 to 5.
  Two curves: Human (amber orange, clear impulse peak at $t{=}0$ reaching $\approx$3.2\,mg,
  decaying smoothly to baseline by $\sim$60\,ms) and LLM Agent (navy blue, flat at noise
  floor $\approx$0.2\,mg throughout).
  Shaded region $[0, 60]$\,ms in light amber. Error band ($\pm$1 std) around each curve.
  Vertical dashed line at $t{=}0$ labeled ``Touch Event.''
  White background, unified palette.]}
  \caption{Event-locked accelerometer response to touch events.
  Real finger contact (human) produces a detectable mechanical impulse ($\approx$3.2\,mg,
  SNR $\approx$8.4\,dB, decay $\lesssim$60\,ms);
  ADB-injected touch events (agent) produce no accelerometer response,
  confirming the structural physical causality gap.}
  \label{fig:obs1}
\end{figure}

\subsection{Observation 2: Touch Interval Distributions Are Separable}
\label{sec:obs2}

Fig.~\ref{fig:obs2} shows the empirical probability density of inter-touch intervals
$\Delta t_i = \tau_{i+1} - \tau_i$ for human and agent sessions.%
Human intervals follow a lognormal distribution
(mean 423\,ms, CV 0.68), consistent with neuromotor
variability in voluntary motor tasks~\cite{fitts1954}.%
AutoGLM and AppAgent sessions, which invoke a cloud vision-language model between each
action, exhibit a bimodal distribution: one peak at $\approx$130\,ms
corresponding to rapid repeat-tap actions, and a second peak at $\approx$1{,}050\,ms
corresponding to LLM inference plus network round-trip latency.%
ADB-script sessions without LLM inference show a near-uniform distribution with a
characteristic minimum inter-event interval floor of $\approx$47\,ms
imposed by the Android input dispatch scheduler.%

These distributional differences are complementary to the causality signal: an adversary
can adjust touch timing to mimic the human lognormal distribution (A2 adversary,
Section~\ref{sec:adversary}), but cannot simultaneously produce a physical impulse at
each artificially timed event.%

\begin{figure}[t]
  \centering
  \textcolor{green}{[Figure: Probability density of inter-touch intervals $\Delta t$.
  X-axis: interval duration (ms), log scale, 50 to 5000, ticks at 50, 100, 200, 500, 1000, 2000, 5000.
  Y-axis: probability density.
  Three curves: Human (amber orange, lognormal shape, single mode around 400--600\,ms),
  LLM Agent cloud (navy blue, bimodal: small peak at $\sim$130\,ms and larger peak at
  $\sim$1{,}050\,ms), ADB script (steel gray, near-uniform with hard floor at $\sim$47\,ms).
  Two vertical dashed navy lines marking the two LLM-agent peaks.
  White background, unified palette.]}
  \caption{Inter-touch interval distributions.
  Human intervals follow a lognormal distribution (mean 423\,ms, CV 0.68) from neuromotor
  variability; LLM agents show a bimodal signature from inference latency; scripted ADB
  shows a near-uniform distribution with a scheduler-imposed floor at $\approx$47\,ms.}
  \label{fig:obs2}
\end{figure}

\subsection{Observation 3: Joint Features Form Linearly Separable Clusters}
\label{sec:obs3}

To assess joint discriminability, we extract a two-dimensional feature vector per
session window: (i)~causality score $C = \frac{1}{N_e}\sum_{i=1}^{N_e} \mathbf{1}[|\delta(t_i^+)| > \theta]$,
the fraction of touch events followed by a detectable impulse, and
(ii)~bimodality index $B$ of the inter-touch interval distribution~\cite{bimodality1985}.%
Fig.~\ref{fig:obs3} shows the t-SNE projection~\cite{vandermaaten2008tsne} of
480 human windows and 320 agent windows across all
session types.%
The three classes (human, autonomous agent, overlay agent) form well-separated clusters
(silhouette score $\approx$0.81), confirming that even simple instantiations
of the two primitives yield strong separability.%
This motivates PCCN (Section~\ref{sec:pccn}), which learns richer representations
of both signals end-to-end.%

\begin{figure}[t]
  \centering
  \textcolor{green}{[Figure: t-SNE 2D scatter plot.
  Three clusters: Human (amber orange squares), Autonomous Agent (navy blue circles),
  Overlay Agent (steel gray triangles).
  Three visually separated cluster regions with minimal overlap.
  Human cluster: upper-left region.
  Autonomous Agent cluster: lower-right region.
  Overlay Agent cluster: upper-right region (partially overlapping with Human, but distinct).
  Axes: t-SNE Dim 1 / t-SNE Dim 2. No tick numbers (schematic).
  Faint dashed ellipse outline around each cluster in the corresponding color.
  Legend in top-right corner.
  White background, unified palette.]}
  \caption{t-SNE projection of joint causality score and bimodality index features.
  Human, autonomous agent, and overlay agent windows form three well-separated clusters
  (silhouette score 0.81), validating the discriminability of the two physical primitives
  across all threat classes.}
  \label{fig:obs3}
\end{figure}

\section{System Model and Adversary Model}
\label{sec:sysmodel}

\subsection{System Model}

PhysCausal involves two principals: a \textit{client device} (the smartphone under
test) and a \textit{detection server} operated by a platform or application provider
(e.g., a mobile bank, ad exchange, or app store).%

\noindent\textbf{Client device.}
The client runs a lightweight PhysCausal agent embedded in the host application.%
The agent operates in one of two server-selected modes.%

In \textit{passive mode}, the agent performs two concurrent actions during the
monitoring window: (i)~records \texttt{TYPE\_LINEAR\_ACCELERATION} at
\texttt{SENSOR\_DELAY\_FASTEST} and (ii)~logs touch-event timestamps reported by the
Android input system.%
No interaction is requested from the user and no visible UI change occurs.%

In \textit{active mode}, the server additionally selects a challenge time
$t^* \in [\frac{T}{4}, \frac{3T}{4}]$ \textit{after} the monitoring window begins and
notifies the app.%
The app presents a natural UI element (e.g., a tappable button or confirmation widget)
at $t^*$.%
The agent records the Android-reported touch timestamp and transmits both streams to
the server.%
Crucially, $t^*$ is communicated only after recording starts, so any pre-recorded
accelerometer stream injected by an adversary cannot contain a correctly timed
impulse transient.%

No inference runs on-device; the full PCCN pipeline executes server-side, keeping
the detection model confidential and preventing white-box adaptive attacks.%

\noindent\textbf{Detection server.}
The server receives the encrypted payload, extracts the linear acceleration signal and
touch-event timestamps, and runs PCCN in passive or active mode.%
In passive mode, PCCN classifies using $\delta(t)$ and $\tau(t)$ alone.%
In active mode, PCCN additionally verifies the causal coupling at $t^*$.%
The server returns a ternary verdict (\textit{human}, \textit{autonomous agent},
or \textit{overlay agent}) to the calling service for downstream enforcement.%

\noindent\textbf{Deployment properties.}
The agent requires no dangerous Android permissions.%
Passive monitoring adds zero user friction.%
Active mode presents a single natural UI element indistinguishable from normal app flow.%
Detection is triggered server-side, so a farm operator cannot anticipate or selectively
suppress individual checks.%

\textcolor{green}{[Figure: System architecture diagram. Left column (Client Device):
smartphone icon at top; three stacked blocks connected by arrows:
(1) ``Linear Acceleration Recording'' (gray, waveform icon, zero-permission label),
(2) ``Touch Event Logging'' (gray, tap-finger icon),
(3) ``Active Mode: Challenge at $t^*$'' (gray, dashed border, clock icon, label
``$t^*$ chosen after recording starts'').
Arrow to right: ``Encrypted Payload.''
Right column (Detection Server): three stacked blocks:
(1) ``Baseline Estimation $\rightarrow$ $\delta(t)$'' (blue),
(2) ``PCCN Classifier'' split into ``Passive: $\delta(t) + \tau(t)$'' and
``Active: $\delta(t) + \tau(t)$ at $t^*$'' (blue, dashed divider),
(3) output diamond ``Human / Autonomous Agent / Overlay Agent'' (amber/navy/gray tricolor).
Arrow from output: ``Verdict to App Backend.''
White background, unified color palette.]}

\subsection{Adversary Model}
\label{sec:adversary}

The adversary is an LLM agent operator whose goal is to cause agent-controlled sessions
to be labeled as human-operated.%
We consider four adversary levels of increasing capability.%

\noindent\textbf{A0 (Default --- no evasion).}
The agent dispatches touch events via ADB or \texttt{AccessibilityService} with no
attempt to modify sensor output.%
The accelerometer shows no impulse at any touch-event timestamp.%
This is trivially detected by PhysCausal.%

\noindent\textbf{A1 (Vibration motor).}
The adversary attaches a vibration motor to the phone chassis and triggers it
synchronously with each injected touch event, producing a mechanical signal in the
accelerometer.%
Motor-induced vibrations have a different spectral signature from real finger-impact
impulses: motors produce narrowband sinusoidal oscillation at the motor's driving
frequency, whereas finger impact produces a broadband transient with characteristic
exponential decay.%
In active mode, the vibration motor must additionally fire at the server-selected $t^*$,
which is communicated only after recording begins; the motor controller can comply, but
the spectral and temporal waveform mismatch remains detectable.%

\noindent\textbf{A2 (Timing mimicry).}
The adversary adjusts inter-touch intervals to match the lognormal distribution of
human timing, defeating the timing-distribution stream of PCCN.%
Active mode remains fully effective: mimicking timing does not produce physical impulses.%
Passive mode retains residual detection from the causality stream alone.%

\noindent\textbf{A3 (Human operator).}
The adversary assigns a human to physically hold and tap the phone for each session,
producing authentic physical signals.%
This eliminates the economic advantage of LLM automation (1:1 human-to-session ratio)
and is treated as outside the practical threat model.%

\section{PhysCausal System Design}
\label{sec:method}

\subsection{System Pipeline}
\label{sec:pipeline}

PhysCausal operates as a five-step pipeline executed across client and server.%

\noindent\textbf{Step 1: Session initialization.}
The detection server issues a monitoring request specifying mode (passive or active) and
window length $T$ (default 30\,s).%
In active mode, the server additionally schedules a challenge time
$t^* \in [\frac{T}{4}, \frac{3T}{4}]$ but withholds $t^*$ until recording has started.%

\noindent\textbf{Step 2: Signal acquisition.}
The PhysCausal agent, embedded in the host application, registers two listeners:
(i)~\texttt{TYPE\_LINEAR\_ACCELERATION} via \texttt{SensorManager} at
\texttt{SENSOR\_DELAY\_FASTEST} (sampling rate $f_s \approx 200$\,Hz on current hardware),
and (ii)~the \texttt{View.OnTouchListener} callback, which records every
\texttt{ACTION\_DOWN} timestamp from the Android input system.%
In active mode, the server transmits $t^*$ over the open encrypted channel after
recording starts; the agent presents the challenge UI element at $t^*$ and records
the resulting touch timestamp.%
At window end, the agent transmits the raw sample buffer and touch-event list to the server.%

\noindent\textbf{Step 3: Preprocessing.}
On the server, three preprocessing steps produce the model inputs.%
First, a moving-median filter (window 50\,ms) estimates the slow background motion
$\bar{a}(t)$, yielding the deviation signal
$\delta_{\mathrm{acc}}(t) = \|\mathbf{a}(t) - \bar{\mathbf{a}}(t)\|_2$
that suppresses postural drift while preserving impulse transients.%
Second, the signal is bandpass-filtered (1--50\,Hz, 4th-order Butterworth) to remove
sensor DC offset and aliased noise.%
Third, $\delta_{\mathrm{acc}}(t)$ is z-score normalized per session to remove
device-to-device gain differences.%
The touch-event list is converted to a sparse indicator vector
$\tau(t) \in \{0,1\}^T$ at the same sampling grid.%
Inter-touch intervals $\mathcal{I} = \{\Delta t_i = \tau_{i+1} - \tau_i\}$ are retained as a separate
input stream.%

\noindent\textbf{Step 4: PCCN inference.}
The Physical Causality Consistency Network (Section~\ref{sec:pccn}) classifies the session.%

\noindent\textbf{Step 5: Verdict.}
The server returns a ternary label (\textit{human}, \textit{autonomous agent},
\textit{overlay agent}) to the calling service for downstream enforcement.%

\subsection{Physical Causality Consistency Network (PCCN)}
\label{sec:pccn}

PCCN jointly processes three input streams---the deviation signal, the touch-event
indicator, and the inter-touch interval distribution---through a shared Transformer
backbone with cross-modal attention.%
Fig.~\ref{fig:pccn} shows the full architecture.%

\begin{figure}[t]
  \centering
  \textcolor{green}{[Figure: PCCN architecture.
  Three vertical input lanes on the left:
  Lane 1 (top): ``$\delta_{\mathrm{acc}}(t)$'' waveform $\rightarrow$ ``Patch Embedding''
    (gray block) $\rightarrow$ ``Transformer Encoder ($L$ layers)'' (blue block)
    $\rightarrow$ $E_{\mathrm{acc}} \in \mathbb{R}^{N_p \times d}$.
  Lane 2 (middle): ``$\tau(t)$ Touch Events'' sparse bar $\rightarrow$
    ``Event Encoding + Positional'' (gray block)
    $\rightarrow$ $E_{\mathrm{touch}} \in \mathbb{R}^{N_e \times d}$.
  Cross-modal attention block (center-right, amber border):
    $Q = E_{\mathrm{touch}}$, $K/V = E_{\mathrm{acc}}$
    $\rightarrow$ $E_{\mathrm{causal}} \in \mathbb{R}^{N_e \times d}$
    labeled ``Causal Impulse Detector.''
  Lane 3 (bottom): ``$\mathcal{I} = \{\Delta t_i\}$'' histogram bars $\rightarrow$
    ``Distribution Encoder (LogNormal + BiModal heads)'' (gray block)
    $\rightarrow$ $E_{\mathrm{timing}} \in \mathbb{R}^d$.
  Two output heads on the right:
    Top head: ``Passive Head: GAP($E_{\mathrm{acc}}$) $\rightarrow$ MLP'' $\rightarrow$ 2-class ŷ.
    Bottom head: ``Active Head:
      $[\mathrm{GAP}(E_{\mathrm{acc}}) \| \mathrm{GAP}(E_{\mathrm{causal}}) \| E_{\mathrm{timing}}]$
      $\rightarrow$ MLP'' $\rightarrow$ 3-class ŷ (Human / Autonomous / Overlay).
  White background, unified navy/amber/gray palette.]}
  \caption{PCCN architecture. Three input streams feed a shared Transformer backbone
  with cross-modal attention that detects causal coupling between accelerometer impulses
  and touch events. Passive and active heads share all upstream parameters.}
  \label{fig:pccn}
\end{figure}

\noindent\textbf{Stream 1: Accelerometer encoder.}
The deviation signal $\delta_{\mathrm{acc}}(t) \in \mathbb{R}^T$ is divided into
non-overlapping patches of length $p$ (default $p{=}40$ samples, 200\,ms at 200\,Hz).%
Each patch is projected to a $d$-dimensional embedding via a learned linear layer,
yielding $N_p = T/p$ patch tokens.%
Learnable positional embeddings are added, and the token sequence is processed by
an $L$-layer Transformer encoder with multi-head self-attention ($H$ heads) and
feed-forward width $4d$.%
The output $E_{\mathrm{acc}} \in \mathbb{R}^{N_p \times d}$ encodes temporal structure
across the full monitoring window.%

\noindent\textbf{Stream 2: Touch-event encoder.}
Touch events $\{t_i^e\}_{i=1}^{N_e}$ define $N_e$ event tokens.%
Each token is initialized to a learned event embedding $\mathbf{e}_{\mathrm{evt}}$
plus a continuous sinusoidal positional encoding derived from the event's
absolute timestamp within the window.%
This produces $E_{\mathrm{touch}} \in \mathbb{R}^{N_e \times d}$.%

\noindent\textbf{Cross-modal attention: causal impulse detector.}
A single multi-head cross-attention layer queries the accelerometer stream at each
touch-event position:
\begin{equation}
E_{\mathrm{causal}} = \mathrm{Softmax}\!\left(\frac{Q_{\mathrm{touch}}\,K_{\mathrm{acc}}^\top}{\sqrt{d}}\right) V_{\mathrm{acc}},
\label{eq:crossattn}
\end{equation}
where $Q_{\mathrm{touch}} = E_{\mathrm{touch}} W^Q$,
$K_{\mathrm{acc}} = E_{\mathrm{acc}} W^K$,
$V_{\mathrm{acc}} = E_{\mathrm{acc}} W^V$
are learned projections.%
The attention weights concentrate on the patch tokens near each touch event, and
$E_{\mathrm{causal}} \in \mathbb{R}^{N_e \times d}$ encodes whether a causal
accelerometer impulse is present at each touch position.%
For the overlay agent class, genuine human touches produce impulses across the session,
but the server-selected challenge touch at $t^*$ is executed by ADB and produces no
impulse; the cross-attention at that token position captures this localized absence.%

\noindent\textbf{Stream 3: Timing distribution encoder.}
Inter-touch intervals $\mathcal{I} = \{\Delta t_i\}$ are embedded via two parallel branches:
(i)~a log-transform head computes summary statistics $[\mu_{\log}, \sigma_{\log}, \mathrm{skew}]$
of $\log(\Delta t_i)$, capturing lognormal shape parameters;
(ii)~a bimodality head computes the bimodality coefficient
$B = (\hat{\gamma}^2 + 1) / (\hat{\kappa} + 3\frac{(n-1)^2}{(n-2)(n-3)})$~\cite{bimodality1985}
and the ratio of sub-300\,ms events (characteristic of LLM cloud inference latency).%
Both branch outputs are concatenated and projected to $E_{\mathrm{timing}} \in \mathbb{R}^d$
via an MLP.%
This stream captures the lognormal (human), bimodal-latency (LLM cloud agent), and
near-uniform (ADB script) signatures identified in Section~\ref{sec:obs2}.%

\noindent\textbf{Passive head.}
A global average pooling (GAP) of $E_{\mathrm{acc}}$ followed by a two-layer MLP
produces a 2-class logit (human vs.\ any agent).%
Passive mode uses this head alone; it operates without any touch-event input by pooling
pure accelerometer context.%

\noindent\textbf{Active head.}
The three streams are fused by concatenation:
\begin{equation}
\mathbf{h} = \bigl[\mathrm{GAP}(E_{\mathrm{acc}}) \;\|\; \mathrm{GAP}(E_{\mathrm{causal}}) \;\|\; E_{\mathrm{timing}}\bigr] \in \mathbb{R}^{3d}.
\label{eq:fusion}
\end{equation}
A two-layer MLP with dropout (rate 0.2) maps $\mathbf{h}$ to 3-class logits
(\textit{human}, \textit{autonomous agent}, \textit{overlay agent}).%
The passive head is a subgraph of this: the active head learns to exploit
physical causality evidence at $t^*$ in addition to background accelerometer context
and timing statistics.%

\noindent\textbf{Gated active-mode contribution.}
To share parameters across modes, a scalar gate $g = \sigma(W_g \mathbf{h})$ weights
the causal term at inference: if the session is declared passive (no $t^*$),
the model sets $\mathrm{GAP}(E_{\mathrm{causal}}) = \mathbf{0}$, reducing the
active head to a timing-aware passive classifier.%

\subsection{Three-Stage Training}
\label{sec:training}

\noindent\textbf{Stage 1: Self-supervised masked pre-training.}
We pre-train the Transformer encoder on a large corpus of unlabeled smartphone
accelerometer recordings using masked autoencoding (MAE)~\cite{he2022mae}.%
A random 75\% of patch tokens are masked; the encoder processes visible tokens,
and a lightweight decoder reconstructs the masked patches.%
The pre-training objective is the mean squared error on masked positions:
\begin{equation}
\mathcal{L}_{\mathrm{MAE}} = \frac{1}{|\mathcal{M}|} \sum_{i \in \mathcal{M}} \|\hat{\mathbf{p}}_i - \mathbf{p}_i\|_2^2.
\label{eq:mae}
\end{equation}
This stage teaches the encoder to model the statistical structure of device
motion without any label supervision.%

\noindent\textbf{Stage 2: Supervised contrastive fine-tuning.}
We fine-tune the encoder on labeled sessions using a combined loss:
\begin{equation}
\mathcal{L}_{\mathrm{stage2}} = \mathcal{L}_{\mathrm{CE}} + \lambda_1 \mathcal{L}_{\mathrm{TFC}} + \lambda_2 \mathcal{L}_{\mathrm{SCL}},
\label{eq:stage2}
\end{equation}
where $\mathcal{L}_{\mathrm{CE}}$ is cross-entropy on class labels,
$\mathcal{L}_{\mathrm{TFC}}$ is time-frequency consistency contrastive loss~\cite{tfc2022}
that enforces agreement between time-domain and frequency-domain views of the same
session, and $\mathcal{L}_{\mathrm{SCL}}$ is supervised contrastive loss~\cite{supcon2020}
that pulls same-class embeddings together and pushes different-class embeddings apart.%
Default weights: $\lambda_1 = 0.5$, $\lambda_2 = 0.1$.%

\noindent\textbf{Stage 3: Cross-modal head training.}
With the Transformer encoder frozen, we train the cross-modal attention layer,
the timing distribution encoder, and both classification heads jointly on labeled
touch--accelerometer pairs.%
This stage teaches the model to exploit physical causality coupling without corrupting
the pre-trained temporal features.%
We apply label smoothing (0.1) and use AdamW~\cite{adamw2019} with a cosine learning
rate schedule.%

\subsection{Session Aggregation}
\label{sec:aggregation}

For long monitoring windows ($T > 30$\,s), we segment the session into non-overlapping
windows of 30\,s and run PCCN independently on each segment.%
The session-level verdict is the majority vote of segment-level predictions, with ties
broken by the highest-confidence class.%
This windowed approach allows the system to detect agents that insert human-like delays
in early touches but revert to scripted timing in later interactions.%

\subsection{Complexity and Overhead}
\label{sec:complexity}

PCCN has $\approx$\textcolor{red}{XX}\,M parameters.%
Server-side inference on a 30-s window takes $\approx$\textcolor{red}{XX}\,ms on a
single CPU core, adding negligible latency to the authentication flow.%
The client-side agent uses zero dangerous permissions, consumes
$\approx$\textcolor{red}{XX}\,mAh over a 30-s session (measured on a Pixel 8 at
\texttt{SENSOR\_DELAY\_FASTEST}), and runs in the existing application process without
a dedicated background service.%

\section{Related Work}
\label{sec:related}

\subsection{LLM GUI Agent Security}
\label{sec:rw_agents}

LLM-driven mobile GUI agents have advanced rapidly in capability and deployment.%
AppAgent~\cite{appagent2024} and AutoGLM~\cite{autoglm2024} both achieve
task-completion rates above 36\% on general phone benchmarks, with AutoGLM commercially
deployed.%
CogAgent~\cite{cogagent2024} achieves state-of-the-art GUI navigation using a
dual-resolution 18B VLM without system API access; Mobile-Agent~\cite{mobileagent2024}
enables vision-only phone operation without XML metadata.%
SeeAct~\cite{seeact2024} demonstrates that GPT-4V can act as a generalist web agent when
paired with effective grounding, and AndroidWorld~\cite{androidworld2024} standardizes
evaluation across 116 tasks on 20 real Android apps.%
Anthropic's Computer Use~\cite{anthropiccu2024} extends the paradigm to desktop interfaces.%
The security implications of agent-controlled mobile devices have been documented in the
fraud context: mass account registration, advertising click abuse, welfare farming, and
unauthorized payment confirmation~\cite{llmagentfraud2025,clickfraud2021,adffraud2021}.%
Existing defenses targeting LLM agents focus on prompt injection~\cite{agentinjection2024}
and policy-level access control~\cite{agentpolicy2024}, rather than detecting whether
an active session is agent-controlled.%
Our work addresses this orthogonal detection gap.%

\subsection{Mobile Bot and Click-Farm Detection}
\label{sec:rw_bot}

Classical mobile bot detection relies on device attestation (SafetyNet, Play Integrity API)
and emulator fingerprinting~\cite{emulatordetect2022}.%
Phone farms defeat both: real hardware passes attestation and carries genuine device
identifiers.%
Environment-aware evasion techniques~\cite{sandboxevasion2022} further demonstrate that
distinguishing real devices from sandboxes via user-artifact differences defeats
state-of-the-art mobile sandboxes.%
Statistical behavioral methods---touch trajectory analysis, UI interaction graphs---were
designed for scripted bots~\cite{frauddroid2017} and are bypassed by LLM agents that
adapt to UI state.%
The Android accessibility framework has been extensively misused by malware for
cross-app data theft and unauthorized transaction fraud~\cite{a11yabuse2021,dva2024},
illustrating the breadth of the software injection threat surface.%
Large-scale measurement studies confirm that click-farm-based invalid traffic already
operates at industrial scale~\cite{clickfraud2021,adffraud2021}, motivating detection
primitives that remain valid against real-hardware phone farms.%
BeCAPTCHA~\cite{becaptcha2021} combined accelerometer and touch kinematics for bot
detection in a CAPTCHA-style challenge, but its feature engineering targets
repetitive fixed-script bots and does not model the physical causality gap.%
Turing Test on Screen~\cite{ttos2026} is the closest prior work: it applies kinematic
features to LLM agent detection and reports strong results against vanilla agents, but
explicitly acknowledges that physical sensor spoofing lies outside its scope.%
Our work closes this gap with a structurally unforgeable physical primitive.%

\subsection{Sensor-Based Liveness and User Authentication}
\label{sec:rw_liveness}

Sensor-based liveness detection has been explored in the context of fingerprint and
face recognition~\cite{sensorliveness2023} and touch-screen authentication~\cite{touchauth2020}.%
USENIX Security has shown that even liveness is insufficient: puppet attacks can replay
captured finger motion, requiring behavioral biometrics as a complementary
layer~\cite{livenessnotenough2020}.%
TouchPass~\cite{touchpass2020} authenticates users via chassis vibrations during
touch---a physical signal related to, but distinct from, the causality impulse we exploit.%
Behavioral biometric systems such as MMAuth~\cite{mmauth2022} and
AuthentiSense~\cite{authentisense2023} achieve high accuracy for continuous
authentication, but remain vulnerable to synthesis attacks because they measure
statistical patterns rather than physical necessity.%
Modern CAPTCHA evaluations~\cite{captchastudy2023} confirm that purely behavioral
challenges are increasingly ineffective, with state-of-the-art solvers bypassing
them at high rates.%
BioMotoTouch~\cite{biomotouch2026} characterizes the mechanical impulse response of
real finger contact as measured by the accelerometer, providing the physical model
underlying Observation~1 (Section~\ref{sec:obs1}).%
iStelan~\cite{istelan2023} demonstrates that capacitive touchscreen contact produces
an EM transient detectable by the permission-free magnetometer.%
Zero-permission motion sensors have further been shown to leak speech
content~\cite{accear2022} and private voice-assistant responses~\cite{stealthyimu2023},
and PhyScout~\cite{physcout2024} detects multi-sensor spoofing via spatio-temporal
consistency---illustrating the richness of physical information available at zero
privilege.%
Our work adapts the accelerometer-based physical causality chain to the LLM agent
detection problem, introducing a new three-class threat model (including the novel
overlay agent class) and a cross-modal attention architecture optimized for this setting.%

\subsection{Cross-Modal and Self-Supervised Learning for Sensor Data}
\label{sec:rw_ml}

Self-supervised pre-training on sensor time-series has produced strong generalizable
representations.%
TF-C~\cite{tfc2022} enforces consistency between time-domain and frequency-domain views
via contrastive learning.%
PatchTST~\cite{patchtst2023} segments time series into patch-level tokens for
channel-independent Transformer processing with masked pre-training, directly
motivating our patch-based accelerometer encoder in Stage~1.%
SimMTM~\cite{simmtm2023} recovers masked time points via weighted neighbor aggregation
on the data manifold, providing an alternative masked pre-training objective for
sensor data.%
FOCAL~\cite{focal2023} uses factorized shared-and-private encoders for multi-sensor
cross-modal contrastive pre-training.%
CA-TCC~\cite{catcc2023} extends contrastive learning to semi-supervised sensor HAR.%
TS-TCC~\cite{tstcc2021} applies temporal and contextual contrastive objectives to IMU data.%
Crossformer~\cite{crossformer2023} introduces cross-dimension attention for multivariate
time series, capturing dependencies across sensor channels; our PC-CWA adapts the
cross-attention paradigm to the asymmetric, physics-constrained relationship between
the touch-event and accelerometer streams.%
CausalFormer~\cite{causalformer2024} discovers causal graphs from multivariate time
series via interpretable attention; our PC-CWA differs by encoding \textit{known}
physical causal structure as a hard architectural constraint rather than learning
causality from data.%
Our PCCN Stage~2 training loss builds on TF-C~\cite{tfc2022} and SupCon~\cite{supcon2020},
while our cross-modal attention head is informed by the cross-dimensional processing
principles of FOCAL and Crossformer, adapted to the asymmetric causality structure
of the accelerometer--touch relationship.%

\section{Evaluation}
\label{sec:eval}

\subsection{Data Collection}
\label{sec:data}

\noindent\textbf{Human sessions.}
We recruited \textcolor{red}{XX} participants (age range \textcolor{red}{XX--XX},
\textcolor{red}{XX\%} male) across three naturalistic interaction tasks:
(T1)~social-media browsing (scrolling, tapping posts, entering search queries),
(T2)~mobile banking form filling (entering account numbers, amounts, confirming transactions),
and (T3)~prompted UI tapping (reacting to server-selected challenge buttons at random times).%
Participants held the phone in their preferred grip throughout; no instruction was given
regarding grip posture.%
Each session lasted 60\,s; we collected \textcolor{red}{XX} sessions per participant
across \textcolor{red}{XX} device models spanning four families:
Pixel (6/7/8/8a), Galaxy (S22/S23/A54), OnePlus (11/12), and Xiaomi (12/13)
(\textcolor{red}{XX} models total).%
All data was collected with IRB approval; participants provided informed written consent.%

\noindent\textbf{Agent sessions.}
Agent sessions were generated using three frameworks:
(i)~AutoGLM~\cite{autoglm2024} controlling an identical UI workflow via its default
ADB + cloud-VLM backend;
(ii)~AppAgent~\cite{appagent2024} performing the same tasks;
(iii)~custom ADB scripts replaying touch sequences without LLM inference.%
For each framework, we additionally tested A2-adversary (timing mimicry) sessions in
which the injection timestamps were resampled to match the lognormal distribution of
our human cohort.%
All agent sessions used the phone resting on a flat, non-vibrating surface (A0 default).%
A1-adversary sessions used an off-the-shelf vibration motor (ERM, 13\,mm diameter)
attached to the phone chassis and triggered by the ADB control host.%
A total of \textcolor{red}{XX} agent sessions were collected across frameworks and adversary levels.%

\noindent\textbf{Active-mode sessions.}
For active-mode evaluation, a separate split of \textcolor{red}{XX} human and
\textcolor{red}{XX} agent sessions included server-selected challenge events at
$t^* \in [\frac{T}{4}, \frac{3T}{4}]$ with a live challenge UI element displayed
at $t^*$.%

\subsection{Baselines}
\label{sec:baselines}

We compare PCCN against four baselines.%
\textbf{BeCAPTCHA-Accel}~\cite{becaptcha2021}: the accelerometer-only binary classifier
from BeCAPTCHA, retrained on our dataset.%
\textbf{TToS-Features}~\cite{ttos2026}: touch-kinematic features (coordinate centroid,
velocity profile entropy, inter-touch interval statistics) from Turing Test on Screen,
fed to a gradient-boosted tree classifier.%
\textbf{MiniRocket}~\cite{minirocket2021}: a state-of-the-art time-series classifier
applied to the raw accelerometer signal.%
\textbf{PCCN-Passive}: our model's passive head only (ablation).%

\subsection{Metrics}
\label{sec:metrics}

We report five metrics: overall accuracy (ACC), macro F1, AUC-ROC, equal error rate (EER),
false positive rate (FPR, human misclassified as agent), and false negative rate (FNR,
agent misclassified as human).%
EER is the standard metric for INFOCOM liveness detection papers.%
All results use a 5-fold cross-validation with a zero-overlap device-model split:
no device model in the test fold appears in the training fold.%

\subsection{Main Classification Results}
\label{sec:main}

Table~\ref{tab:main} reports passive-mode and active-mode results under the A0 (no evasion)
adversary condition.%

\begin{table}[t]
\centering
\caption{Classification results under A0 adversary. Passive mode: 2-class (Human vs.\ Agent).
Active mode: 3-class (Human / Autonomous / Overlay). All metrics are mean $\pm$ std
over 5 device-model folds.}
\label{tab:main}
\begin{tabular}{lccccc}
\toprule
Method & Mode & ACC & F1 & AUC & EER \\
\midrule
BeCAPTCHA-Accel~\cite{becaptcha2021} & Passive & \textcolor{red}{XX.X} & \textcolor{red}{XX.X} & \textcolor{red}{XX.X} & \textcolor{red}{XX.X} \\
TToS-Features~\cite{ttos2026}        & Passive & \textcolor{red}{XX.X} & \textcolor{red}{XX.X} & \textcolor{red}{XX.X} & \textcolor{red}{XX.X} \\
MiniRocket~\cite{minirocket2021}     & Passive & \textcolor{red}{XX.X} & \textcolor{red}{XX.X} & \textcolor{red}{XX.X} & \textcolor{red}{XX.X} \\
PCCN-Passive (ours)                  & Passive & \textcolor{red}{XX.X} & \textcolor{red}{XX.X} & \textcolor{red}{XX.X} & \textcolor{red}{XX.X} \\
\midrule
PCCN-Full (ours)                     & Active  & \textcolor{red}{XX.X} & \textcolor{red}{XX.X} & \textcolor{red}{XX.X} & \textcolor{red}{XX.X} \\
\bottomrule
\end{tabular}
\end{table}

\noindent\textbf{Key findings.}
PCCN-Full achieves \textcolor{red}{>XX\%} accuracy in active mode, outperforming all
passive baselines by \textcolor{red}{XX} percentage points.%
The kinematic baseline (TToS-Features) performs near chance on A2-timing-mimicry sessions
(see Section~\ref{sec:adversarial}), confirming the limitation identified in
Section~\ref{sec:bg_kinematic}.%
BeCAPTCHA-Accel achieves moderate accuracy on A0 sessions but degrades substantially
on A1 (vibration motor) sessions; its feature engineering does not capture the spectral
distinction between motor oscillation and finger-impact broadband transients.%

\subsection{Cross-Device Generalization}
\label{sec:crossdevice}

Fig.~\ref{fig:crossdevice} shows per-fold accuracy under the zero-overlap device-model
split, with training and test folds partitioned by manufacturer family.%
PCCN-Full achieves \textcolor{red}{>XX\%} across all four family splits, demonstrating
that the physical causality gap is device-agnostic.%
The smallest accuracy is observed when training on Pixel/OnePlus and testing on Xiaomi,
attributable to the higher baseline accelerometer noise floor on Xiaomi hardware;
PCCN compensates via the per-session z-score normalization in preprocessing.%

\begin{figure}[t]
  \centering
  \textcolor{green}{[Figure: Grouped bar chart. X-axis: test device family
  (Pixel, Galaxy, OnePlus, Xiaomi). Y-axis: accuracy (\%), 80--100 range.
  Four bar groups, each with five bars for: BeCAPTCHA-Accel (steel gray),
  TToS-Features (light blue), MiniRocket (green), PCCN-Passive (amber orange),
  PCCN-Full (navy blue). PCCN-Full bar consistently tallest.
  Error bars ($\pm$1 std). Dashed horizontal reference line at human-level benchmark.
  White background, unified palette.]}
  \caption{Cross-device generalization under a zero-overlap device-model split.
  PCCN-Full maintains $>$\textcolor{red}{XX\%} accuracy across all four device
  families, confirming device-agnostic operation.}
  \label{fig:crossdevice}
\end{figure}

\subsection{Adversarial Robustness}
\label{sec:adversarial}

Table~\ref{tab:adversarial} reports active-mode accuracy of PCCN-Full and baselines
under A0--A2 adversary conditions.%

\begin{table}[t]
\centering
\caption{Active-mode accuracy under adversary levels A0--A2.
Best result per adversary level in bold.}
\label{tab:adversarial}
\begin{tabular}{lcccc}
\toprule
Method & A0 (default) & A1 (vibration) & A2 (timing mimicry) \\
\midrule
BeCAPTCHA-Accel  & \textcolor{red}{XX.X} & \textcolor{red}{XX.X} & \textcolor{red}{XX.X} \\
TToS-Features    & \textcolor{red}{XX.X} & \textcolor{red}{XX.X} & \textcolor{red}{XX.X} \\
MiniRocket        & \textcolor{red}{XX.X} & \textcolor{red}{XX.X} & \textcolor{red}{XX.X} \\
PCCN-Full (ours) & \textbf{\textcolor{red}{XX.X}} & \textbf{\textcolor{red}{XX.X}} & \textbf{\textcolor{red}{XX.X}} \\
\bottomrule
\end{tabular}
\end{table}

\noindent\textbf{A1 (vibration motor).}
PCCN-Full retains \textcolor{red}{>XX\%} accuracy against A1 adversaries.%
The cross-modal attention layer learns to distinguish broadband finger-impact transients
from the narrowband sinusoidal oscillation produced by the ERM motor.%
The spectral centroid of the motor signal is concentrated at the motor's driving frequency
(\textcolor{red}{$\approx$XX}\,Hz), while genuine finger impact spans 1--50\,Hz with
an exponential spectral rolloff.%
The Transformer encoder captures this distinction via learned frequency-sensitive filters
in its self-attention heads.%

\noindent\textbf{A2 (timing mimicry).}
Timing mimicry reduces TToS-Features accuracy to near-chance but has no effect on
PCCN-Full, which relies on physical impulse presence rather than timing statistics.%
PCCN-Passive (which uses accelerometer context only, without timing stream) similarly
retains high accuracy under A2, confirming that the physical causality signal drives detection.%

\subsection{Overlay Agent Detection}
\label{sec:overlay}

Fig.~\ref{fig:overlay} shows the detection window sweep for overlay agent sessions as
a function of the number of challenge events $t^*$ observed in the session.%
With a single challenge event (\textit{n=1}), PCCN-Full achieves
\textcolor{red}{>XX\%} detection accuracy at FPR \textcolor{red}{$<$XX\%}.%
The key insight is that an overlay agent replaces exactly the challenged touch with an ADB
injection; the cross-modal attention head detects the impulse absence at that specific
token while genuine human impulses remain at all other positions.%

\begin{figure}[t]
  \centering
  \textcolor{green}{[Figure: Line plot. X-axis: number of challenge events $n$ (1, 2, 3, 4, 5).
  Y-axis: detection accuracy (\%), 60--100.
  Two lines: PCCN-Full (navy blue, rising from $\sim$XX\% at n=1 to $\sim$XX\% at n=5),
  PCCN-Passive (amber orange, flat near $\sim$XX\% regardless of n, since it cannot
  exploit active challenges).
  Shaded confidence bands. Dashed horizontal line at 95\% reference.
  White background, unified palette.]}
  \caption{Overlay agent detection accuracy vs.\ number of server-selected challenge
  events. A single challenge event is sufficient for $>$\textcolor{red}{XX\%} detection;
  passive mode cannot distinguish overlay agents from human sessions.}
  \label{fig:overlay}
\end{figure}

\subsection{Ablation Study}
\label{sec:ablation}

Table~\ref{tab:ablation} ablates the four components of PCCN-Full: the accelerometer
Transformer encoder, cross-modal attention (causal impulse detector), timing distribution
encoder, and Stage~1 MAE pre-training.%

\begin{table}[t]
\centering
\caption{Ablation study. Active-mode accuracy under A0. Each row removes one component.}
\label{tab:ablation}
\begin{tabular}{lcc}
\toprule
Configuration & ACC & EER \\
\midrule
PCCN-Full                             & \textcolor{red}{XX.X} & \textcolor{red}{XX.X} \\
w/o Cross-modal attention             & \textcolor{red}{XX.X} & \textcolor{red}{XX.X} \\
w/o Timing distribution encoder      & \textcolor{red}{XX.X} & \textcolor{red}{XX.X} \\
w/o MAE pre-training (Stage~1)        & \textcolor{red}{XX.X} & \textcolor{red}{XX.X} \\
Passive head only (no active mode)    & \textcolor{red}{XX.X} & \textcolor{red}{XX.X} \\
\bottomrule
\end{tabular}
\end{table}

Removing cross-modal attention causes the largest drop (\textcolor{red}{XX} pp),
confirming that explicit causal impulse detection is the primary signal.%
Removing the timing encoder causes a smaller drop (\textcolor{red}{XX} pp) that is
concentrated in the A2-timing-mimicry condition (where timing mimicry eliminates
the timing signal, leaving physical causality as the sole discriminant).%
MAE pre-training contributes \textcolor{red}{XX} pp, consistent with its role in
generalizing the encoder to unseen device models.%

\subsection{Attention Visualization}
\label{sec:attention}

Fig.~\ref{fig:attention} shows the cross-modal attention weights averaged over
\textcolor{red}{XX} human and \textcolor{red}{XX} agent test sessions, plotted
as a heatmap of attention weight vs.\ time offset relative to each touch event.%
Without explicit supervision on the [0,\,80]\,ms window, PCCN independently
concentrates attention at offsets 0--60\,ms post-touch for human sessions.%
Agent sessions show uniformly diffuse attention, consistent with the absence of
a localized impulse.%
This post-hoc interpretability confirms that PCCN has learned the physical causality
mechanism, not a spurious correlation.%

\begin{figure}[t]
  \centering
  \textcolor{green}{[Figure: Two-panel heatmap. Left panel (Human sessions):
  X-axis: time offset relative to touch event (ms), $-100$ to $+200$.
  Y-axis: touch event index (ordered by session time).
  Color: cross-modal attention weight (white$\rightarrow$amber, higher = more attention).
  Strong amber band at 0--60\,ms, consistent across all event rows.
  Right panel (Agent sessions): same axes.
  Color: uniformly pale (near-zero attention everywhere).
  Vertical dashed line at $t{=}0$ in both panels.
  Shared color bar on right.
  White background, unified palette.]}
  \caption{Cross-modal attention visualization. For human sessions (left), PCCN
  concentrates attention in the [0,\,60]\,ms post-touch window without explicit supervision,
  rediscovering the physical impulse response. Agent sessions (right) show diffuse attention,
  reflecting the absence of a causal impulse.}
  \label{fig:attention}
\end{figure}

\subsection{System Overhead}
\label{sec:overhead}

We measure client-side overhead on a Pixel 8 running Android 14.%
Sensor registration at \texttt{SENSOR\_DELAY\_FASTEST} and touch-event logging consume
\textcolor{red}{XX}\,mAh over a 30-s session (measured via Battery Historian),
equivalent to \textcolor{red}{XX\%} of idle power consumption.%
No background service is required; the agent runs within the host application process.%
Server-side PCCN inference on a 30-s window takes \textcolor{red}{XX}\,ms on a
single Intel Xeon core, well within the latency budget of any authentication flow.%
Peak memory consumption on the server is \textcolor{red}{XX}\,MB per concurrent session.%

\section{Conclusion}
\label{sec:conclusion}

LLM GUI agents defeat device attestation and kinematic detection because they run on
real hardware and can synthesize statistically plausible touch trajectories.%
The distinguishing property of an agent-controlled session is not its behavioral
statistics but its \textit{physical origin}: software touch injection never actuates
the touchscreen hardware, so no mechanical impulse propagates through the chassis
and no accelerometer transient is produced.%

PhysCausal exploits this structural impossibility.%
A zero-permission linear accelerometer records the impulse signature of genuine finger
contact; the Physical Causality Consistency Network jointly processes the accelerometer
stream and touch-event timestamps via cross-modal attention, detecting the presence or
absence of causal impulse coupling at each touch position.%
A server-selected challenge time $t^*$, communicated only after recording begins,
makes pre-recorded impulse injection infeasible and enables detection of the novel
overlay agent threat---where an LLM agent injects fraudulent taps into a live human
session.%

We train PCCN in three stages---masked accelerometer pre-training, supervised
contrastive fine-tuning, and cross-modal head training on touch--accelerometer
pairs---achieving \textcolor{red}{>XX\%} accuracy under a zero-overlap device-model
split across \textcolor{red}{XX} Android device families.%
Post-hoc attention visualization confirms that PCCN independently rediscovers the
[0,\,60]\,ms physical impulse window, providing an interpretable and
physics-grounded explanation for every detection decision.%

The broader implication is that \textit{physical causality signals} constitute an
unforgeable primitive for mobile session liveness: their presence requires genuine
human contact that no software injection can fabricate, and their targeted absence
at a server-specified challenge time cannot be concealed by any pre-recorded stream.%
Future work includes multi-sensor fusion (combining accelerometer and magnetometer
primitives) when elevated permissions are available, federated deployment preserving
user privacy while aggregating detection evidence across installations, and
investigation of A3-level physical tethering countermeasures to bound the economic
cost imposed on sophisticated farm operators.%


\bibliographystyle{IEEEtran}
\bibliography{ref}

\end{document}
